\section{Smart Contracts for Public Issue Management and Opinion Polling}

In governance systems, to improve the transparency, efficiency and citizen participation, smart contracts have been explored as the good mechanism. Smart contracts can automate the governance rules and incentives by immutable code. It can support models like civic intelligence governance (CIG). In CIG, citizen engagement is encouraged through on chain participation and token based incentives\cite{MDPI_2023_CivicIntelligence}. Literature shows that approaches like CIG can improve participation rates and knowledge sharing while reducing relances on centralized intermediaries. In the context of multi organizational governance smart contracts also been used within dual blockchain architectures. This method separates transaction data from contract logic. It balance the transparency and privacy in Government to Government and Business to Business processes\cite{IEEE_2021_SmartCon}.

% Smart contracts also help  business process reengineering in governance processes by directly embedding policy rules into the executable logic. Also the smart contracts help to do the tendering processes and public procurement, automate bid submission, evaluation, and contract execution by making these processes with zero chances for corruption and also increasing the accountability \cite{IEEE_2021_SmartCon}, \cite{GovernanceNow_2023_SmartContractsToolkit}, \cite{AICERTs_2023_FutureOfGovernance}. Also Similar benefits can be seen in blockchain based land registry systems by enabling immutable property records, automated ownership transfers and enforcement of land use policies. So again this reduces the
% fraud and inefficiencies in administratives \cite{ResearchGate_2020_LandRegistry}, \cite{MDPI_2023_NamibianLandRegistration}.

Smart contracts also support business processes in governance by embedding policy rules directly into executable logic. Bid submission, evaluation and contract execution are some of the use cases in the public procurement and tendering processes that smart contract hold its value. It also reduce the oppertunities for corruption and increase the accountability\cite{IEEE_2021_SmartCon}, \cite{GovernanceNow_2023_SmartContractsToolkit}, \cite{AICERTs_2023_FutureOfGovernance}. Similar benefits are observed in the blockchain based land registry systems. Smart contracts have enable immutable property records automated ownership transfers and enforcement of land use policies. This has reduced fraud and administrative inefficiencies\cite{ResearchGate_2020_LandRegistry}, \cite{MDPI_2023_NamibianLandRegistration}.

% If we consider voting and opinion polling, blockchain based systems make sure the security and transparency is protected by ensuring immutability and verifiability of votes while preserving voter anonymity \cite{ResearchGate_2024_ScalableEVoting}, \cite{IEEE_2024_BPVot}. Advanced mechanisms like Quadratic Voting and blockchain enabled participatory budgeting increases the citizen involvement in decision making. They allow more visibility in preference representation and transparent allocation of public resources \cite{ResearchGate_2021_QuadraticVoting}, \cite{UPSpace_2022_ParticipatoryBudgeting}. Eventhough there are many advantages like that, existing studies also highlight some challenges related to scalability, privacy, and system complexity. Particularly when deploying voting and polling solutions for large and diverse populations \cite{ResearchGate_2024_ScalableEVoting}, \cite{JournalOfCybersecurity_2020_BlockchainVoting}.

In the context of voting and opinion polling, smart contracts offer enhanced security, transparency, immutability and verifiability inherited from blockchain\cite{ResearchGate_2024_ScalableEVoting}, \cite{IEEE_2024_BPVot}. There are advance mechanisms such as Quadratic Voting and also blockchain enabled participatory budgeting. Those mechanisms further extend citizen involvement in decision making. It have expressive preference representation and transparent allocation of public resources\cite{ResearchGate_2021_QuadraticVoting}, \cite{UPSpace_2022_ParticipatoryBudgeting}. In addition to these advantages, there exists some challenges related to scalability, privacy and system complexity when deploying voting and polling systems to large and diverse populations\cite{ResearchGate_2024_ScalableEVoting}, \cite{JournalOfCybersecurity_2020_BlockchainVoting}.

